\section{Artin-Wedderburn}

\begin{theorem}
    Let
    \[
        {}_RV = V_1 \oplus \cdots \oplus V_n
    \]
    be a direct sum of $R$-modules, and suppose all $V_i$ are isomorphic to each other.
    Then
    \[
        \operatorname{End}_R({}_RV) \cong M_n\bigl(\operatorname{End}_R(V_i)\bigr).
    \]
\end{theorem}

\begin{proof}
    For each $i$, let
    \[
        \iota_i:V_i\hookrightarrow V
        \qquad\text{and}\qquad
        \pi_i:V\to V_i
    \]
    be the canonical inclusion and projection. Fix isomorphisms
    \[
        \theta_i:V_1\xrightarrow{\sim} V_i
        \qquad (i=1,\dots,n).
    \]

    For $\sigma\in \operatorname{End}_R(V)$, define
    \[
        \varphi_{ij}(\sigma):=\theta_i^{-1}\pi_i\,\sigma\,\iota_j\theta_j
        \in \operatorname{End}_R(V_1).
    \]
    Thus we obtain a map
    \[
        \Phi:\operatorname{End}_R(V)\longrightarrow M_n(E_1),
        \qquad
        \sigma\longmapsto \bigl(\varphi_{ij}(\sigma)\bigr),
    \]
    where
    \[
        E_1:=\operatorname{End}_R(V_1).
    \]

    We first check that $\Phi$ is a ring homomorphism.

    Additivity is immediate. For multiplicativity, let $\sigma,\beta\in \operatorname{End}_R(V)$. Since
    \[
        1_V=\sum_{j=1}^n \iota_j\pi_j,
    \]
    we have
    \begin{align*}
        \varphi_{ik}(\beta\sigma)
         & =\theta_i^{-1}\pi_i\,\beta\sigma\,\iota_k\theta_k             \\
         & =\theta_i^{-1}\pi_i\,\beta
        \left(\sum_{j=1}^n \iota_j\pi_j\right)
        \sigma\,\iota_k\theta_k                                          \\
         & =\sum_{j=1}^n
        \theta_i^{-1}\pi_i\,\beta\,\iota_j\pi_j\,\sigma\,\iota_k\theta_k \\
         & =\sum_{j=1}^n
        \bigl(\theta_i^{-1}\pi_i\,\beta\,\iota_j\theta_j\bigr)
        \bigl(\theta_j^{-1}\pi_j\,\sigma\,\iota_k\theta_k\bigr)          \\
         & =\sum_{j=1}^n \varphi_{ij}(\beta)\varphi_{jk}(\sigma).
    \end{align*}
    Hence
    \[
        \Phi(\beta\sigma)=\Phi(\beta)\Phi(\sigma).
    \]
    So $\Phi$ is a ring homomorphism.

    Next we show that $\Phi$ is injective. Suppose $\Phi(\sigma)=0$. Then
    \[
        \varphi_{ij}(\sigma)=0
        \qquad\text{for all } i,j,
    \]
    so
    \[
        \theta_i^{-1}\pi_i\,\sigma\,\iota_j\theta_j=0
        \qquad\text{for all } i,j.
    \]
    Since each $\theta_i$ and $\theta_j$ is an isomorphism, this implies
    \[
        \pi_i\,\sigma\,\iota_j=0
        \qquad\text{for all } i,j.
    \]
    Therefore $\pi_i(\sigma(v))=0$ for every $v\in V_j$ and every $i,j$. Hence
    \[
        \sigma(V_j)=0
        \qquad\text{for all } j.
    \]
    Because
    \[
        V=V_1\oplus\cdots\oplus V_n,
    \]
    it follows that $\sigma=0$. Thus $\Phi$ is injective.

    Finally we prove surjectivity. Let
    \[
        (a_{ij})\in M_n(E_1).
    \]
    Define
    \[
        \sigma:=\sum_{i,j=1}^n \iota_i\theta_i\,a_{ij}\,\theta_j^{-1}\pi_j
        \in \operatorname{End}_R(V).
    \]
    Then for each $i,j$,
    \begin{align*}
        \varphi_{ij}(\sigma)
         & =\theta_i^{-1}\pi_i
        \left(\sum_{r,s=1}^n \iota_r\theta_r\,a_{rs}\,\theta_s^{-1}\pi_s\right)
        \iota_j\theta_j        \\
         & =\sum_{r,s=1}^n
        \theta_i^{-1}(\pi_i\iota_r)\theta_r\,a_{rs}\,\theta_s^{-1}(\pi_s\iota_j)\theta_j.
    \end{align*}
    Now
    \[
        \pi_i\iota_r=
        \begin{cases}
            1_{V_i}, & i=r,     \\
            0,       & i\neq r,
        \end{cases}
        \qquad
        \pi_s\iota_j=
        \begin{cases}
            1_{V_j}, & s=j,     \\
            0,       & s\neq j,
        \end{cases}
    \]
    so all terms vanish except the one with $r=i$ and $s=j$. Hence
    \[
        \varphi_{ij}(\sigma)
        =\theta_i^{-1}\theta_i\,a_{ij}\,\theta_j^{-1}\theta_j
        =a_{ij}.
    \]
    Therefore
    \[
        \Phi(\sigma)=(a_{ij}),
    \]
    and $\Phi$ is surjective.

    Thus $\Phi$ is a ring isomorphism, and we conclude that
    \[
        \operatorname{End}_R(V)\cong M_n\bigl(\operatorname{End}_R(V_1)\bigr).
    \]
\end{proof}

\begin{definition}
    A ring $R$ is called semisimple if ${}_R R$ is semisimple as a left $R$-module.
\end{definition}

\begin{lemma}
    For a unital ring $R$, the following are equivalent:
    (1) $R$ is semisimple.
    (2) Every left $R$-module is semisimple.
\end{lemma}

\begin{proof}
    If ${}_R R$ is semisimple, then every free module is semisimple, and every module is a quotient of a free module.
    Use closure properties of semisimple modules to deduce that every left $R$-module is semisimple.
    The converse is immediate by applying the statement to ${}_R R$ itself.
\end{proof}

\begin{lemma}
    For an $R$-module ${}_RM$, there is a standard identification
    \[
        \operatorname{Hom}_R({}_RR_R, {}_RM) \cong {}_RM,
        \qquad
        f \longmapsto f(1).
    \]
\end{lemma}

\begin{proof}
    Define
    \[
        \Phi:\operatorname{Hom}_R({}_RR_R,{}_RM)\longrightarrow {}_RM,
        \qquad
        \Phi(f)=f(1).
    \]

    We first show that $\Phi$ is bijective.

    For any $f\in \operatorname{Hom}_R({}_RR_R,{}_RM)$ and any $r\in R$, since $f$ is
    left $R$-linear, we have
    \[
        f(r)=f(r\cdot 1)=r\,f(1).
    \]
    Thus $f$ is completely determined by the element $f(1)\in M$.

    Conversely, for each $m\in M$, define
    \[
        \Psi(m):{}_RR_R\longrightarrow {}_RM,
        \qquad
        \Psi(m)(r)=rm
        \quad (r\in R).
    \]
    Then $\Psi(m)$ is left $R$-linear, because for all $a,r\in R$,
    \[
        \Psi(m)(ar)=(ar)m=a(rm)=a\,\Psi(m)(r).
    \]
    Hence $\Psi(m)\in \operatorname{Hom}_R({}_RR_R,{}_RM)$.

    Now for $m\in M$,
    \[
        (\Phi\circ \Psi)(m)=\Phi(\Psi(m))=\Psi(m)(1)=1m=m.
    \]
    And for $f\in \operatorname{Hom}_R({}_RR_R,{}_RM)$ and $r\in R$,
    \[
        (\Psi\circ\Phi)(f)(r)=\Psi(f(1))(r)=r\,f(1)=f(r).
    \]
    So $\Psi\circ\Phi=\mathrm{id}$ and $\Phi\circ\Psi=\mathrm{id}$, hence $\Phi$ is a bijection.

    Finally, we check that this identification is $R$-linear. Recall that
    $\operatorname{Hom}_R({}_RR_R,{}_RM)$ carries the left $R$-module structure
    \[
        (a\cdot f)(r)=f(ra)
        \qquad (a,r\in R).
    \]
    Then for $a\in R$ and $f\in \operatorname{Hom}_R({}_RR_R,{}_RM)$,
    \[
        \Phi(a\cdot f)=(a\cdot f)(1)=f(1a)=f(a)=a\,f(1)=a\,\Phi(f).
    \]
    Therefore $\Phi$ is an isomorphism of left $R$-modules:
    \[
        \operatorname{Hom}_R({}_RR_R,{}_RM)\cong {}_RM.
    \]
\end{proof}

\begin{corollary}
    For any ring $R$, there is a ring isomorphism
    \[
        \operatorname{End}_R({}_R R)\cong R^{\operatorname{op}},
        \qquad
        f\longmapsto f(1).
    \]
\end{corollary}

\begin{proof}
    Consider the map
    \[
        \Phi:\operatorname{End}_R({}_R R)\longrightarrow R,
        \qquad
        \Phi(f)=f(1).
    \]
    By the preceding lemma with $M=R$, this map is a bijection of sets.

    We now check the ring structure. Let $f_1,f_2\in \operatorname{End}_R({}_R R)$. Then
    \[
        \Phi(f_1\circ f_2)
        =(f_1\circ f_2)(1)
        =f_1\bigl(f_2(1)\bigr).
    \]
    Since $f_1$ is left $R$-linear, for any $r\in R$ we have
    \[
        f_1(r)=f_1(r\cdot 1)=r\,f_1(1).
    \]
    Hence
    \[
        f_1\bigl(f_2(1)\bigr)=f_2(1)\,f_1(1).
    \]
    Therefore
    \[
        \Phi(f_1\circ f_2)=\Phi(f_2)\Phi(f_1),
    \]
    where the product on the right is the multiplication in $R$. This means exactly that
    \[
        \Phi(f_1\circ f_2)=\Phi(f_1)\cdot_{\operatorname{op}}\Phi(f_2)
    \]
    when the codomain is viewed as $R^{\operatorname{op}}$.

    Thus $\Phi$ is a ring isomorphism
    \[
        \operatorname{End}_R({}_R R)\cong R^{\operatorname{op}}.
    \]
\end{proof}

\begin{lemma}
    For any ring $R$ and any $n\ge 1$, there is a ring isomorphism
    \[
        M_n(R)^{\operatorname{op}}\cong M_n\bigl(R^{\operatorname{op}}\bigr),
        \qquad
        A\longmapsto A^{\top}.
    \]
\end{lemma}

\begin{proof}
    Define
    \[
        \Psi:M_n(R)^{\operatorname{op}}\longrightarrow M_n\bigl(R^{\operatorname{op}}\bigr),
        \qquad
        \Psi(A)=A^{\top}.
    \]
    This map is clearly additive and bijective, since transpose is its own inverse.

    It remains to check multiplicativity. Let
    \[
        A=(a_{ij}),\qquad B=(b_{ij})\in M_n(R).
    \]
    Recall that multiplication in $M_n(R)^{\operatorname{op}}$ is reversed, so
    \[
        A\cdot_{\operatorname{op}} B = BA
    \]
    with the usual matrix product on the right-hand side.

    Now fix $i,j$. Then
    \[
        \bigl(\Psi(A\cdot_{\operatorname{op}} B)\bigr)_{ij}
        =\bigl((BA)^{\top}\bigr)_{ij}
        =(BA)_{ji}
        =\sum_{k=1}^n b_{jk}a_{ki}.
    \]
    On the other hand, in the ring $R^{\operatorname{op}}$ the product is reversed, so
    \[
        \bigl(\Psi(A)\Psi(B)\bigr)_{ij}
        =\sum_{k=1}^n (A^{\top})_{ik}\cdot_{R^{\operatorname{op}}}(B^{\top})_{kj}
        =\sum_{k=1}^n a_{ki}\cdot_{R^{\operatorname{op}}} b_{jk}
        =\sum_{k=1}^n b_{jk}a_{ki}.
    \]
    Thus
    \[
        \Psi(A\cdot_{\operatorname{op}} B)=\Psi(A)\Psi(B).
    \]
    So $\Psi$ is a ring homomorphism, hence a ring isomorphism.

    Therefore
    \[
        M_n(R)^{\operatorname{op}}\cong M_n\bigl(R^{\operatorname{op}}\bigr).
    \]
\end{proof}

\begin{theorem}
    (Artin-Wedderburn) Let $A$ be a semisimple Artinian ring.
    Then
    \[
        A \cong M_{n_1}(D_1) \oplus \cdots \oplus M_{n_r}(D_r),
    \]
    where each $D_i$ is a division ring.
    More concretely, if
    \[
        {}_A A \cong S_1^{\oplus n_1} \oplus \cdots \oplus S_r^{\oplus n_r}
    \]
    with pairwise nonisomorphic simple $A$-modules $S_i$, then
    \[
        D_i = \operatorname{End}_A(S_i)^{op}.
    \]
\end{theorem}

\begin{proof}
    Decompose $\operatorname{End}_A(A)$ using the decomposition of the regular module.
    Use
    \[
        \operatorname{Hom}_A(S_i,S_j)=0
        \quad (i \neq j)
    \]
    and Schur's lemma to identify the diagonal blocks with division rings.
    Then use the previous matrix-ring result to rewrite
    \[
        \operatorname{End}_A\bigl(S_i^{\oplus n_i}\bigr)
    \]
    as
    \[
        M_{n_i}\bigl(\operatorname{End}_A(S_i)\bigr)=M_{n_i}(D_i^{op}).
    \]
    Therefore
    \[
        A \cong \operatorname{End}_A(A)^{op} \cong M_{n_1}(D_1) \oplus \cdots \oplus M_{n_r}(D_r).
    \]
\end{proof}



\begin{definition}
    Let $R$ be a commutative ring with $1$, and let $A$ be an associative ring with $1$.
    An $R$-algebra is an $R$-module $A$ whose multiplication is compatible with the $R$-action.
    Equivalently, there is a ring homomorphism
    \[
        R \to A
    \]
    whose image lies in the center $Z(A)$.
\end{definition}

\begin{remark}
    If $R=k$ is a field and $A$ is finite-dimensional over $k$, then $A$ is called a finite-dimensional $k$-algebra.
    Such an algebra is Noetherian and Artinian as a ring/module object.
\end{remark}

\begin{example}
    Let $R$ be a commutative ring with $1$. The following are $R$-algebras:
    \begin{enumerate}
        \item $R[x_1,\dots,x_n]$.
        \item $M_n(R)$.
        \item Let $G$ be a finite group. The group algebra
              \[
                  RG=\left\{\sum_{g\in G} a_g g \;\middle|\; a_g\in R\right\}
              \]
              is an $R$-algebra, with multiplication defined by
              \[
                  \left(\sum_{g\in G} a_g g\right)\left(\sum_{h\in G} b_h h\right)
                  =
                  \sum_{r\in G}\left(\sum_{\substack{g,h\in G\\ gh=r}} a_g b_h\right) r.
              \]
    \end{enumerate}
\end{example}

\begin{lemma}
    Let $A$ be a $k$-algebra, where $k$ is a field, and let $S$ be a simple $A$-module.
    Then $\operatorname{End}_A(S)$ is a division ring, and it is naturally a $k$-algebra.
\end{lemma}

\begin{proof}
    We first show that $\operatorname{End}_A(S)$ is a division ring.

    Let
    \[
        0\neq f\in \operatorname{End}_A(S).
    \]
    Since $f$ is $A$-linear, both $\ker(f)$ and $\operatorname{im}(f)$ are $A$-submodules of $S$.
    Because $S$ is simple and $f\neq 0$, we have $\operatorname{im}(f)\neq 0$, hence
    \[
        \operatorname{im}(f)=S.
    \]
    Also $\ker(f)\neq S$ since $f\neq 0$, so by simplicity,
    \[
        \ker(f)=0.
    \]
    Thus $f$ is bijective. Since the inverse of a bijective $A$-linear map is again $A$-linear,
    $f^{-1}\in \operatorname{End}_A(S)$. Therefore every nonzero element of
    $\operatorname{End}_A(S)$ is invertible, so $\operatorname{End}_A(S)$ is a division ring.

    Next we show that $\operatorname{End}_A(S)$ is naturally a $k$-algebra.

    Because $A$ is a $k$-algebra, $S$ becomes a $k$-vector space via
    \[
        \lambda\cdot s := (\lambda 1_A)s
        \qquad (\lambda\in k,\ s\in S).
    \]
    Now let $f\in \operatorname{End}_A(S)$. Then for $\lambda\in k$ and $s\in S$,
    \[
        f(\lambda s)
        =
        f\bigl((\lambda 1_A)s\bigr)
        =
        (\lambda 1_A)f(s)
        =
        \lambda f(s).
    \]
    Hence every $A$-linear endomorphism of $S$ is automatically $k$-linear. So
    \[
        \operatorname{End}_A(S)\subseteq \operatorname{End}_k(S)
    \]
    as a $k$-subspace. Equivalently, there is a $k$-algebra homomorphism
    \[
        k\longrightarrow \operatorname{End}_A(S),
        \qquad
        \lambda\longmapsto \lambda\,\operatorname{id}_S,
    \]
    whose image lies in the center of $\operatorname{End}_A(S)$. Therefore
    $\operatorname{End}_A(S)$ is naturally a $k$-algebra.
\end{proof}

\begin{lemma}
    If $A$ is finite-dimensional over $k$, then $\operatorname{End}_A(S)$ is finite-dimensional over $k$.
\end{lemma}

\begin{proof}
    Choose $0\neq s\in S$. Since $S$ is simple, the cyclic submodule $As$ is nonzero, hence
    \[
        As=S.
    \]
    Consider the map
    \[
        \varphi:A\longrightarrow S,
        \qquad
        a\longmapsto as.
    \]
    This is a surjective $k$-linear map. Therefore
    \[
        \dim_k S\leq \dim_k A<\infty.
    \]
    So $S$ is finite-dimensional over $k$.

    Since every $A$-endomorphism of $S$ is $k$-linear, we have
    \[
        \operatorname{End}_A(S)\subseteq \operatorname{End}_k(S)
    \]
    as $k$-vector spaces. Hence
    \[
        \dim_k \operatorname{End}_A(S)
        \leq
        \dim_k \operatorname{End}_k(S)
        =
        (\dim_k S)^2
        <\infty.
    \]
    Thus $\operatorname{End}_A(S)$ is finite-dimensional over $k$.
\end{proof}

\begin{lemma}
    If $A$ is finite-dimensional over $k$ and $k$ is algebraically closed, then
    \[
        \operatorname{End}_A(S)=k.
    \]
\end{lemma}

\begin{proof}
    By the previous lemma, $S$ is finite-dimensional over $k$, and
    $\operatorname{End}_A(S)$ is a finite-dimensional $k$-algebra.

    Let
    \[
        f\in \operatorname{End}_A(S).
    \]
    Since $S$ is a finite-dimensional $k$-vector space and $k$ is algebraically closed,
    the $k$-linear operator $f:S\to S$ has an eigenvalue $\lambda\in k$. Hence
    \[
        f-\lambda \operatorname{id}_S
    \]
    has a nonzero kernel, so it is not injective.

    But $f-\lambda \operatorname{id}_S$ is still an element of $\operatorname{End}_A(S)$.
    Since $\operatorname{End}_A(S)$ is a division ring by the first lemma, every nonzero
    element is invertible, hence injective. Therefore
    \[
        f-\lambda \operatorname{id}_S=0.
    \]
    So
    \[
        f=\lambda \operatorname{id}_S.
    \]
    This shows that every $A$-endomorphism of $S$ is scalar. Under the natural embedding
    \[
        k\hookrightarrow \operatorname{End}_A(S),
        \qquad
        \lambda\mapsto \lambda\operatorname{id}_S,
    \]
    we therefore obtain
    \[
        \operatorname{End}_A(S)=k.
    \]
\end{proof}

\begin{theorem}
    If $k$ is algebraically closed and $A$ is a finite-dimensional semisimple $k$-algebra, then
    \[
        A \cong M_{n_1}(k) \oplus \cdots \oplus M_{n_r}(k).
    \]
\end{theorem}

\begin{corollary}
    Let $A$ be a finite-dimensional semisimple $k$-algebra, where $k$ is a field.

    In any decomposition
    \[
        {}_A A \cong S_1^{n_1} \oplus \cdots \oplus S_r^{n_r},
    \]
    where the $S_i$'s are pairwise non-isomorphic simple $A$-modules, we have that
    $S_1,\dots,S_r$ form a complete set of representatives of the isomorphism classes
    of simple $A$-modules.

    If moreover $k=\overline{k}$, then $n_i=\dim_k S_i$ and
    \[
        \dim_k A = n_1^2+\cdots+n_r^2.
    \]
\end{corollary}

\begin{proof}
    Let $S$ be a simple module. Then
    \[
        S \cong A/M
    \]
    where $M$ is a maximal submodule.

    By semisimplicity,
    \[
        {}_A A = M \oplus S.
    \]
    Hence all simple submodules appear in $\{S_1,\dots,S_r\}$.

    Since $k=\overline{k}$, we have
    \[
        A \cong M_{n_1}(k)\oplus \cdots \oplus M_{n_r}(k).
    \]

    Moreover,
    \[
        M_{n_i}(k)=
        \underbrace{
            \left\{
            \begin{pmatrix}
                *      & \cdots & 0      \\
                \vdots & \ddots & \vdots \\
                *      & \cdots & 0
            \end{pmatrix}
            \right\}
            \oplus \cdots \oplus
            \left\{
            \begin{pmatrix}
                0      & \cdots & *      \\
                \vdots & \ddots & \vdots \\
                0      & \cdots & *
            \end{pmatrix}
            \right\}
        }_{\text{each of them is a simple } M_{n_i}(k)\text{-module}.}
    \]
\end{proof}

